\chapter{Bevezetés}

Az információs technológiák rohamos fejlődése alapjaiban alakította át a kiberbiztonság területét. 
Az új technológiák megjelenésével ugyan jelentősen bővült a védekezés eszköztára, ez azonban korántsem jelenti a fenyegetések megszűnését vagy csökkenését. 
A digitális térben zajló „küzdelem” inkább egy folyamatosan változó, dinamikus folyamatként írható le, 
ahol a támadási és védekezési módszerek egymással párhuzamosan fejlődnek. 
Az internet széles körű elterjedése egy korábbi jelentős fordulópontot jelentett a kiberbiztonság történetében, átalakítva a küzdelem színterét és módját, 
míg napjainkban a mesterséges intelligencia új eszközöket biztosít mind a támadók, mind a védekezők számára.

A köztudatban élő kép, miszerint a kiberbiztonság valós idejű „programozói párbajként” értelmezhető, nem tükrözi a valóság összetettségét. 
A gyakorlatban a hangsúly sokkal inkább a már működő rendszerekben található sérülékenységek feltárásán és kezelésén van. 
Egy informatikai rendszer biztonsága nem csupán a forráskód minőségétől függ, hanem számos egyéb tényező együttes hatásának eredménye. 
Ilyen tényező például a titkosítás hiánya, a nem megfelelő jogosultságkezelés, a hibás konfigurációk, és nem utolsó sorban az emberi tényező is.

Az egyik leggyakoribb támadási forma az \textbf{adathalászat} (phishing). Az ilyen típusú támadások során a támadók megtévesztő, 
gyakran hitelesnek tűnő üzenetek segítségével próbálnak érzékeny adatokat megszerezni a felhasználóktól. Az adathalász e-mailek különösen veszélyesek, 
mivel kihasználják a felhasználók figyelmetlenségét, illetve döntéshozatali korlátait. Mindez rámutat arra, 
hogy a kizárólag felhasználói döntésekre építő védekezés nem tekinthető elegendőnek.

Ebből következően a modern kiberbiztonsági megoldások egyik kulcseleme az automatizált védelem alkalmazása, mely a valóságban a természetes nyelv terén, 
illetve tapasztalat alapján egyéb területeken is inkább nevezhető félautomatának. 
Az e-mail szolgáltatók fejlett szűrőrendszereket használnak annak érdekében, hogy a nem kívánt vagy kártékony 
üzeneteket még a felhasználóhoz való eljutás előtt azonosítsák és elkülönítsék habár ez nem jelenti, hogy nem kerül kéretlen üzenet a levélfiókba. 
Lehetőséget adnak a felhasználóknak, hogy megtekintsék és esetleg vissza állítsák a tévesen nem kívántank címkézett e-maileket. 

Ezt a folyamatot klasszifikációnak, osztályba sorolásnak tekintjük. Ez lehet egy, kettő vagy akár több osztályú probléma is. 
Gépi módszerek számos fajtája alkalmazható, és ugyancsak sokféle adatból lehet döntést hozni egy e-mail kéretlen voltáról. 
A szöveg alapú megközelítések mellett a metaadatokból, hálózati jellemzőkből és viselkedési mintákból automata vagy félautomata módon információt 
kinyerő modellek és ezek konkrét megvalósítása végtelen számú megoldást nyújthat a feladatra.

A beérkező üzenetek nagy mennyisége és folyamatosan növekvő változatossága miatt a modern rendszerek jellemzően 
\textbf{gépi tanulási} (Machine Learning) módszerekre épülnek, amelyek képesek alkalmazkodni az újonnan megjelenő mintázatokhoz, 
valamint a dinamikusan változó és egyre kifinomultabb támadási technikákhoz. E megközelítések közül különösen kiemelkednek a 
\textbf{mélytanulási} (Deep Learning, DL) módszerek, amelyek a hagyományos gépi tanulási eljárásokhoz képest számos jelentős előnnyel rendelkeznek.

A mélytanulási modellek egyik legfontosabb tulajdonsága, hogy jelentősen csökkentik a kézi \textbf{jellemzőtervezés} (Feature Engineering) szükségességét, 
amely a klasszikus módszerek esetében gyakran időigényes, szakértelmet igénylő és számításigényes folyamat. 
Ezzel szemben a mély neurális hálózatok képesek közvetlenül a nyers adatokból automatikusan kinyerni a releváns jellemzőket és mintázatokat, 
ezáltal a tanulási folyamat sokkal inkább adatvezéreltté és kevésbé emberi beavatkozást igénylővé válik. 
Ennek eredményeként a modell nemcsak hatékonyabban képes kezelni a nagy mennyiségű és komplex adatokat, hanem jobb általánosító képességet is mutathat új, 
korábban nem látott példák esetén.

Az egyik megközelítés a \textbf{kéretlen levél} (spam) és \textbf{valódi levél} (ham) kategóriák osztályba sorolására a természetes nyelvű szekvenciák vizsgálata. 
Ezen dolgozat ezt veszi alapul, és vizsgálja \textbf{mélytanulási} (Deep Learning) módszerekkel. 
Időben vagy térben folytonos sorrendű adatok feldolgozására már korlátozottabb az eszköztár. 
A természetes nyelvi szekvenciák feldolgozására különösen alkalmasak a \textbf{hosszú-rövid távú memória} (Long Short-Term Memory, LSTM) 
architektúrán alapuló visszacsatolt neurális hálózatok, amelyek a nyelvben előforduló statisztikai mintázatokat tanulják meg. 
A szöveges e-maileket sorozatjellegű adatokként kezelik, ahol az egyes szavak és nyelvi mintázatok jelentése nagyban függ az őket megelőző kontextustól. 
A neurális hálózat a korábban látott szókapcsolatok és statisztikai összefüggések alapján tanulja meg felismerni azokat a mintákat, 
amelyek jellemzőek lehetnek a \textbf{kéretlen} (spam) vagy a \textbf{valódi} (ham) levelekre. 
Az osztályozási döntés során nem kizárólag egyes kulcsszavak jelenlétét vizsgálja, hanem azok sorrendjét, gyakoriságát és környezetét is figyelembe veszi. 
Ennek köszönhetően a modell képes összetettebb nyelvi és stilisztikai mintázatok felismerésére is, amelyek hozzájárulnak az e-mailek pontosabb osztályozásához.

A legmodernebb nyelvi modellek működési elve alapvetően továbbra is szekvenciák modellezésére épül, azonban a felépítésük nagyban eltér a korábbi módszerektől. 
A \textbf{nagy nyelvi modellek} (Large Language Model, LLM) jellemzően a transformer architektúrán alapulnak, és előtanítás által általános "nyelvi tudást" szereznek. 
A tanító korpusz hatalmas mennyíségű természetes nyelvi adat melyből nyelvi reprezentációkat és statisztikai összefüggéseket sajátítanak el az adott nyelvről. 
Ennek köszönhetően zajtűrőbbek, kevesebb előfeldolgozást igényelnek tanítás és \textbf{következtetés} (inference) közben.

A dolgozat eredeti célja a mélytanulási módszerek közül a hosszú-rövid távú memória (Long Short-Term Memory, LSTM) architektúra 
és a modern nagy nyelvi modellek (Large Language Model, LLM) vizsgálata, különös tekintettel azok szövegfeldolgozási és osztályozási képességeire. 
A LSTM-alapú megközelítés a dolgozat során gyakorlati formában, implementációval és kísérletekkel kerül bemutatásra, beleértve a modell felépítését, 
tanítását és kiértékelését. Ezzel szemben a nagy nyelvi modellek említésre kerülnek a következtetések levonásakor, a különbségek kiemelésével.

A kutatás további célja a modell megfelelő felépítésének és optimális paraméterezésének kialakítása mellett 
az előfeldolgozási lépések részletes és szisztematikus meghatározása is, mivel ezek alapvetően befolyásolják a tanítás hatékonyságát és a végső teljesítményt. 
Kiemelt szempont a modell pontosságának és egyéb teljesítménymutatóinak átfogó vizualizálása és elemzése, 
amely lehetővé teszi a különböző konfigurációk közötti összehasonlítást és a tanulási folyamat viselkedésének jobb megértését. 
Emellett a dolgozat egyik fontos vizsgálati iránya annak feltárása is, hogy mekkora mennyiségű tanítóadat szükséges egy stabil, 
megbízhatóan általánosító modell eléréséhez, különös tekintettel arra, hogyan változik a teljesítmény a növekvő adatmennyiség függvényében.

